\documentclass{article}
\usepackage{booktabs}
\usepackage{tabularx}
\usepackage[margin=1in]{geometry}
\begin{document}

\begin{table}[tbp] \centering
\newcolumntype{R}{>{\raggedleft\arraybackslash}X}
\newcolumntype{L}{>{\raggedright\arraybackslash}X}
\newcolumntype{C}{>{\centering\arraybackslash}X}

\caption{Gelbach (2016) decomposition: which controls drive the Simpson sign-flip?}
\label{tab:gelbach}
\begin{tabularx}{\linewidth}{@{}lCCC@{}}

\toprule
&{(1)}&{(2)}&{(3)} \tabularnewline \midrule
{Contribution to vineyard coef. change (b1base - b1full)}&{Delta}&{(SE)}&{p} \tabularnewline
\midrule \addlinespace[\belowrulesep]
Language (french\\_share)&--573.9**&(285.0)&0.044
Religion (catholic\\_share)&--89.1&(92.6)&0.336
Total (sum)&--663.0**&(287.8)&0.021
\bottomrule \addlinespace[\belowrulesep]

\end{tabularx}
\\ \parbox{\linewidth}{\footnotesize Notes: Gelbach (2016) conditional decomposition of the change in the vineyard\_per\_cap coefficient between the bivariate spec (b1base = -178.6) and the KEY spec adding french\_share + catholic\_share (b1full =  484.4). Delta is the contribution of each x2 group to b1base minus b1full (= -663.0). The Simpson sign-flip means deltas are NEGATIVE (the controls move the coefficient UP from negative to positive). Implementation: official b1x2 package by Gelbach (2014), v4.1.0, vendored at libraries/stata/b/. Hand-validated against the b1x2 identity (assertion in 05\_expansion.do sec 10.9). Robust SEs (HC1; b1x2 does not support HC3). Methods reference: analysis/documentation/methods/gelbach\_decomposition.md. Significance: * p<0.10, ** p<0.05, *** p<0.01.}
\end{table}
\end{document}
